\documentclass{article}
\usepackage[T1]{fontenc}
\usepackage{lmodern}
\usepackage{graphicx}
\usepackage{amsmath,amssymb}
\usepackage[a4paper,margin=2cm]{geometry}
\usepackage[absolute,overlay]{textpos}
\setlength{\TPHorizModule}{1mm}
\setlength{\TPVertModule}{1mm}
\usepackage[indonesian]{babel}
\usepackage{hyperref}
\setlength{\emergencystretch}{2em}
% unit: Halaman 58

\begin{document}

% === Halaman 58 (python/native) ===

• Di dalam beberapa cipher blok, permutasi direalisasikan dengan menggunakan 

matriks permutasi yang dinamakan P-box. Entri di dalam P-box menyatakan bit 

dari posisi sebelumnya.

• Contoh P-box di dalam DES:

• Beberapa skema permutasi juga melakukan kompresi (mengurangi bit-bit input) 

atau ekspansi (memperbanyak bit-bit) pada operasi permutasi.


58 

50 

42 

34 

26 

18 

10 

2 

60 

52 

44 

36 

28 

20 

12 

4 

62 

54 

46 

38 

30 

22 

14 

6 

64 

56 

48 

40 

32 

24 

16 

8 

57 

49 

41 

33 

25 

17 

9 

1 

59 

51 

43 

35 

27 

19 

11 

3 

61 

53 

45 

37 

29 

21 

13 

5 

63 

55 

47 

39 

31 

23 

15 

7 

 Artinya, bit ke 58 dari bit-bit input dipindahkan ke posisi pertama, bit ke-50 pada posisi kedua, dst

1

2

3

4

5

6

7

8

11

12

13

14

15

16

17

18

9

10

21

22

23

24

25

26

27

28

19

20

31

32

29

30

1

2

3

4

5

6

7

8

11

12

13

14

15

16

17

18

9

10

21

22

23

24

25

26

27

28

19

20

31

32

29

30

28

29

24

25

21

20

16

17

13

12

4

5

8

9

32

1

Right Half i-1

58

\end{document}
